% =====================================================================================
% Chapter 7 · What to control for
%
% Built from notebooks/05_what_to_control_for.ipynb (body and the LaLonde section §7.10).
% NOT ONE NUMBER IN THIS FILE IS TYPED BY HAND: every figure in the prose is a macro from
% build/macros.tex, generated by cmp.report from the executed notebook. The only literals
% below are (i) the coefficients of the data-generating model in §7.2, which are
% *definitions* rather than results, and (ii) the decision-rule threshold (0.9) and the
% conventional significance level (0.05), which are *conventions* stated in the text.
% =====================================================================================
\chapter{What to control for}
\label{ch:ctrl}

\begin{quote}\small
\emph{An email campaign is about to go to \nbFiveCampaignN{} customers at a fully loaded
\eur{\nbFiveCost} per contact, and the analyst has a customer table with every column the CRM
ever wrote. Regressing spend on the email and nothing else returns \eur{\nbFiveNaiveEst} per
customer; regressing it on the email and the one variable the causal graph demands returns
\eur{\nbFiveAdjEst}, against a planted truth of \eur{\nbFiveTrueAte}. Budgeting on the first
books \eur{\nbFivePhantom} of incremental revenue that does not exist. And the instinct that
looks like caution --- control for everything, since more information cannot hurt --- returns
\eur{\nbFiveKsBar}, puts the truth outside its 94\,\% credible interval, and reverses the
campaign decision from \nbFiveDecision{} to \nbFiveDecisionKs. Its interval is not wider as a
penalty for being wrong: it is \eur{\nbFiveKsBayesWidth} against the correct model's
\eur{\nbFiveBayesWidth}. No diagnostic in the output --- not the standard error, not $\hat R$,
not the posterior predictive --- distinguishes the confident lie from the answer. Only the
graph does.}
\end{quote}

% =====================================================================================
\section{The decision}
\label{sec:ctrl:decision}

Every estimate in every other chapter of this book depends on a choice made \emph{before} any
model runs: which variables to hold fixed. The choice is made in a moment, usually without
argument, and it decides whether the number that emerges is the causal effect or an artefact.
This chapter is about that moment.

The concrete decision is the smallest one in the book. A retailer sends a discount email; the
question is what the email adds to a customer's spend, and whether that exceeds the
\eur{\nbFiveCost} it costs to send, target and fulfil. At \nbFiveCampaignN{} customers the
campaign is worth roughly a quarter of a million euros in contact cost, so the per-customer
effect is multiplied by fifty thousand before anyone sees it: an error of one euro per customer
is an error of \nbFiveCampaignN{} euros in the budget.

The instinct almost every analyst brings to it --- and it is worth stating baldly, because it is
so nearly universal --- is \emph{control for everything you have}. The reasoning feels
unimpeachable. More information cannot hurt. Leaving a variable out risks confounding; putting
it in risks, at worst, a slightly wider interval. Adjustment is insurance, and insurance is
cheap.

Every clause of that reasoning is false, and this chapter demonstrates each falsehood with a
number. Controlling for the wrong variable does not widen the interval; it moves the estimate,
and leaves the interval where it was. Adjusting for a variable measured \emph{before} treatment
can create bias where none existed. Filtering a dataset --- an act no model summary even
records --- is the same operation as adjustment and carries the same danger. And the estimator,
however sophisticated, cannot tell you any of this: the same OLS, the same sampler, the same
immaculate convergence diagnostics come back whether the adjustment set is right or wrong. The
only object that knows the difference is the causal graph, and it must be drawn before the data
is touched.

The chapter therefore inverts the book's usual emphasis. Elsewhere the estimator is the subject
and identification is the preamble. Here identification is the entire subject, and the estimator
is deliberately the plainest one available --- ordinary least squares --- so that nothing can be
credited to it.

% =====================================================================================
\section{The data-generating model}
\label{sec:ctrl:dgm}

A control choice cannot be graded on real data, because on real data the right answer is exactly
what is missing. So the choice is first graded against a world whose answer is known by
construction, and only then trusted on a world whose answer is not. \Cref{sec:ctrl:real}
performs the second half of that bargain on the LaLonde job-training data, where an experimental
benchmark exists; this section builds the first.

The simulator draws \nbFiveN{} independent customers. Each carries a loyalty score
$L \sim \mathcal N(0,1)$ --- observed, recorded in the CRM --- and an \emph{unobserved}
engagement trait $G \sim \mathcal N(0,1)$, the kind of latent disposition no warehouse holds.
Writing $\sigma(z) = 1/(1+e^{-z})$ for the logistic function, the world is
%
\begin{align}
  T = \text{email} &\;\sim\; \text{Bernoulli}\big(\sigma(0.8\,L)\big),
      \label{eq:ctrl:treat}\\
  Y = \text{spend} &\;=\; 20 \;+\; 5\,L \;+\; 6\,T \;+\; 3\,G \;+\; \varepsilon,
      \qquad \varepsilon \sim \mathcal N(0,\,4^{2}), \label{eq:ctrl:outcome}\\
  \text{opened\_email} &\;\sim\; \text{Bernoulli}\big(\sigma(1.5\,T + 0.7\,G)\big),
      \label{eq:ctrl:opened}\\
  \text{responded} &\;\sim\; \text{Bernoulli}\big(\sigma(1.2\,T + 0.05\,(Y - \bar Y))\big).
      \label{eq:ctrl:responded}
\end{align}
%
The planted average treatment effect is the coefficient on $T$ in \cref{eq:ctrl:outcome}:
\eur{\nbFiveTrueAte} of incremental spend per email. That is the grade every estimate in this
chapter is marked against, and it is the last quantity in the chapter that is known for certain.

Four features of this world are deliberate, and each corresponds to a variable a real marketing
table would contain.

\begin{description}[leftmargin=1.1em, style=nextline, font=\normalfont\itshape]
  \item[Loyalty is a confounder.] $L$ appears in \emph{both} \cref{eq:ctrl:treat} and
  \cref{eq:ctrl:outcome}: loyal customers are more likely to be emailed \emph{and} spend more
  anyway. That is the backdoor $T \leftarrow L \rightarrow Y$, and it must be closed.
  \item[\code{responded} is a collider.] It is caused by the email and by spend
  (\cref{eq:ctrl:responded}) and appears nowhere in the spend equation. It cannot possibly
  improve an estimate of $T \to Y$; it can only manufacture bias.
  \item[\code{opened\_email} is the subtler trap.] It is caused by the email and by the
  \emph{unobserved} $G$, which also drives spend. It is therefore a collider too --- but a
  collider on a path through a latent variable, which means a DAG drawn from the observed
  columns alone \emph{cannot see it}. \Cref{sec:ctrl:valid} catches it anyway, from data.
  \item[$G$ is unobserved.] It exists so that the observed graph is incomplete, which is the
  normal condition of applied work. Pretending otherwise is what makes DAG demonstrations
  unconvincing to practitioners.
\end{description}

Two auxiliary worlds, of \nbFiveMiniN{} customers each, isolate the two structures that usually
survive only as pictures. In the first, the email is \emph{randomized} by a coin flip --- so
there is no confounding at all --- and it lifts spend both directly and through site visits,
$T \to M \to Y$: a mediator. In the second, a pre-treatment covariate (last year's coupon
redemptions) is the common effect of two latent traits, one of which drives targeting and the
other of which drives spend: the M-bias structure of \cref{fig:ctrl:dags}, panel four. Both are
priced in \cref{sec:ctrl:wrong}.

\input{build/floats/nb05_dags}

% =====================================================================================
\section{Identification}
\label{sec:ctrl:ident}

\subsection*{The estimand}

In the potential-outcome notation of \citet{rubin1974}, write $Y_i(1)$ for what customer $i$
would spend if emailed and $Y_i(0)$ for what the same customer would spend if not. Exactly one of
the two is ever observed. The quantity the campaign is bought against is the average treatment
effect,
%
\begin{equation}
  \tau \;=\; \E\big[\,Y_i(1) - Y_i(0)\,\big],
  \label{eq:ctrl:estimand}
\end{equation}
%
which \cref{eq:ctrl:outcome} plants at \eur{\nbFiveTrueAte}. Everything in this chapter --- five
regressions, two posteriors, one sensitivity analysis and the euro decision --- targets that
single scalar. What differs between them is not the estimator. It is which variables are held
fixed while it runs.

\subsection*{To control for a variable is to condition on it}

\emph{Controlling for} $W$ --- equivalently \emph{adjusting for}, \emph{conditioning on} --- means
comparing treated and untreated customers who share the same value of $W$, and then averaging
those within-stratum comparisons back together. Including $W$ as a regressor is a linear shortcut
for exactly that operation. So is stratifying. So is matching. And so --- a point
\cref{sec:ctrl:wrong} will make expensive --- is \emph{filtering the dataset on $W$}, which is
conditioning on a single value of it.

Whether the operation helps or hurts depends entirely on the causal role $W$ plays between $T$
and $Y$, and there are four roles (\cref{fig:ctrl:dags}). A \emph{confounder} is a common cause
of both. A \emph{mediator} lies on the causal path. A \emph{collider} is a common effect. An
\emph{M-bias collider} is a common effect of two latent causes, one feeding treatment and the
other feeding the outcome. The rule for the first is \emph{adjust}; for the other three,
\emph{do not} \citep{cinelli2022}.

\subsection*{What blocks a path: d-separation}

The formal machinery is \citeauthor{pearl2009}'s, and it is worth stating as a rule that can be
run by hand, because the software that automates it does nothing more. A \emph{path} between $T$
and $Y$ is any chain of edges connecting them, arrows pointing either way. Given a candidate
control set $Z$, a path is \textbf{blocked} if and only if it contains at least one of
%
\begin{equation}
  \underbrace{A \to B \to C \quad\text{or}\quad A \leftarrow B \to C,
              \qquad\text{with } B \in Z}_{\text{chain or fork: conditioning on the middle node closes it}}
  \label{eq:ctrl:dsep1}
\end{equation}
\begin{equation}
  \underbrace{A \to B \leftarrow C, \qquad\text{with } B \notin Z
              \text{ and no descendant of } B \text{ in } Z}_{%
              \text{collider: closed by default --- conditioning on it \emph{opens} the path}}
  \label{eq:ctrl:dsep2}
\end{equation}
%
where a \emph{descendant} of $B$ is any node reachable from $B$ by following arrows forward.
\Cref{eq:ctrl:dsep1} is the intuitive half: compare customers at the \emph{same} loyalty and the
``loyal people get emailed and spend more anyway'' signal disappears. \Cref{eq:ctrl:dsep2} is the
half that costs money, and it runs the other way. A collider is closed \emph{until you touch it}.
Conditioning on it creates an association that was never there.

The mechanism is \emph{explaining away}, and it deserves a sentence in plain language because
every counter-intuitive result in this chapter reduces to it. Among the customers who responded,
observing a high spend makes the email a \emph{less necessary} explanation for the response ---
something already accounts for it. So within the responder stratum, email and spend read as
negatively related even when they are independent in the population. Conditioning on a common
effect makes its causes compete \citep{elwert2014}.

\subsection*{The backdoor criterion}

\begin{proposition}[Backdoor criterion; \citealp{pearl1995, pearl2009}]
\label{prop:ctrl:backdoor}
A set $Z$ is an \emph{admissible adjustment set} for the effect of $T$ on $Y$ if
%
\begin{enumerate}[label=(\alph*), leftmargin=2.2em, itemsep=.15em]
  \item no element of $Z$ is a descendant of $T$, and
  \item $Z$ blocks every path from $T$ to $Y$ that starts with an arrow \emph{into} $T$.
\end{enumerate}
%
Such paths are the \emph{backdoor paths} --- the routes along which ``who got treated'' leaks
information about the outcome. When (a) and (b) hold, the interventional distribution is
recovered from observational data by the \emph{adjustment formula}
%
\begin{equation}
  \Prob\big(Y \mid do(T{=}t)\big) \;=\; \sum_{z} \Prob\big(Y \mid T{=}t,\, Z{=}z\big)\,\Prob(z),
  \label{eq:ctrl:adjust}
\end{equation}
%
where $do(T{=}t)$ denotes setting $T$ by intervention --- everybody gets the email --- as opposed
to observing $T{=}t$, which the loyal select into.
\end{proposition}

\Cref{eq:ctrl:adjust} reads: \emph{compare treated with untreated inside look-alike strata of
$Z$, then average the strata back with population weights}. Every regression adjustment in this
chapter is a linear shortcut for it.

Running \cref{prop:ctrl:backdoor} by hand on the graph of \cref{sec:ctrl:dgm} takes four lines.
Enumerate the paths from \code{email} to \code{spend}:

\begin{itemize}[leftmargin=*, itemsep=.15em]
  \item $\text{email} \to \text{spend}$ --- the causal path being estimated. Leave it alone;
        blocking it is what a mediator does.
  \item $\text{email} \leftarrow \text{loyalty} \to \text{spend}$ --- starts with an arrow
        \emph{into} \code{email}, hence a backdoor, and open by default at the fork. Any
        admissible $Z$ must contain \code{loyalty}.
  \item $\text{email} \to \text{responded} \leftarrow \text{spend}$ --- a collider at
        \code{responded}, hence already blocked. Putting \code{responded} into $Z$ would open it,
        and \code{responded} is a descendant of $T$, so it fails clause (a) as well.
  \item $\text{email} \to \text{opened\_email}$ --- in the \emph{declared} graph, a dead end: no
        arrow onward to \code{spend}, so it lies on no $T$--$Y$ path at all. Remember this one.
        It is the blind spot \cref{sec:ctrl:valid} exposes.
\end{itemize}

So $Z = \{\text{loyalty}\}$ is admissible, and minimal. The \code{pathmc} enumeration in the
notebook returns exactly \nbFiveNAdjSets{} admissible set, and it is that one; it also warns, as
it must, that a candidate set containing \code{responded} conditions on a collider. Every
identification claim in this chapter is that enumeration, and nothing else.

\subsection*{The assumptions}

\begin{assumption}[Conditional ignorability / no open backdoor]\label{as:ctrl:backdoor}
$\{Y(0), Y(1)\} \perp\!\!\!\perp T \mid Z$: within strata of $Z$, treatment is as good as
randomly assigned. This is what \cref{prop:ctrl:backdoor} delivers when the graph is right, and
it is \emph{not} testable from the data --- no residual plot, no goodness-of-fit statistic and no
convergence diagnostic can see an open backdoor. It is discharged by the graph, and the graph is
an assumption. \Cref{as:ctrl:graph} is where it is put at risk.
\end{assumption}

\begin{assumption}[Positivity / overlap]\label{as:ctrl:positivity}
$0 < \Prob(T{=}1 \mid Z{=}z) < 1$ for every stratum $z$ that carries mass. \Cref{eq:ctrl:adjust}
averages only over strata in which \emph{both} arms exist; a stratum with no treated customers
contributes nothing but an extrapolation. Here it holds by construction, since
$\sigma(0.8L) \in (0,1)$ for every finite loyalty score. On real data, with a high-dimensional
$Z$, it fails silently, and the standard check is a propensity-score histogram by arm, as the
uplift chapter's Step~4 performs it.
\end{assumption}

\begin{assumption}[No adjustment for descendants of $T$]\label{as:ctrl:descendants}
Clause (a) of \cref{prop:ctrl:backdoor}. Mediators, colliders and any post-treatment variable
are excluded. This one \emph{is} checkable --- against the timeline of the data pipeline, not
against the data --- and it is the assumption most often violated in practice, because the
post-treatment columns (\code{opened}, \code{clicked}, \code{responded}) are the ones a
marketing table is richest in.
\end{assumption}

\begin{assumption}[The graph is right]\label{as:ctrl:graph}
The declared DAG encodes the true causal structure: no missing edge, no reversed arrow, no
omitted latent common cause. This is the load-bearing assumption of the whole method, and it is
the one \citet{cinelli2022} call the price of admission. It is not fully testable --- but it is
not \emph{un}testable either, and the difference matters. A DAG implies a list of conditional
independences that must hold in the data. Those are checkable, and \cref{sec:ctrl:valid} checks
them. \textbf{The graph declared here fails one of them.}
\end{assumption}

\begin{assumption}[SUTVA]\label{as:ctrl:sutva}
One version of the treatment, and no interference: one customer's email does not change another
customer's spend \citep{imbens2015}. Household sharing and word of mouth violate it. No test in
this chapter can see it; it is discharged by knowing how the campaign was fielded, and it is
stated here rather than dressed in a diagnostic's clothes.
\end{assumption}

\begin{remark}
\Cref{eq:ctrl:adjust} is nonparametric: it asks for strata, not for a linear model. Every fit in
this chapter uses OLS, which additionally assumes the adjustment enters linearly. That
assumption is true by construction here (\cref{eq:ctrl:outcome} is linear in $L$), and it is a
genuine, separate risk on real data. It is worth separating the two failures in one's head,
because they are repaired by different things: a wrong \emph{functional form} is repaired by a
better estimator, and a wrong \emph{adjustment set} is not repaired by anything.
\end{remark}

% =====================================================================================
\section{The classical baseline}
\label{sec:ctrl:classical}

The discipline of this book is to ask what a competent analyst produces without a sampler. Here
that question is not a detour: it is the whole argument. The estimator is the plainest one there
is --- spend regressed on the email plus an adjustment set $Z$,
%
\begin{equation}
  \text{spend}_i \;=\; \beta_0 \;+\; \tau\,\text{email}_i \;+\; \gamma^{\top} Z_i
                      \;+\; \varepsilon_i ,
  \label{eq:ctrl:ols}
\end{equation}
%
with $\hat\tau$, the coefficient on \code{email}, as the estimate of \cref{eq:ctrl:estimand}. No
likelihood, no prior, no chain; a point estimate and an interval, in milliseconds. It is run
twice, and \emph{the only thing that changes between the two runs is $Z$}.

\subsection*{The covariance choice, argued and then tested}

An interval built on the wrong standard error is a confident lie, so the covariance assumption is
argued from the design rather than defaulted. The rows here are \nbFiveN{} \emph{independent
customers, each measured once}. There is no panel: each customer contributes exactly one row, so
there is no within-cluster correlation for a cluster-robust estimator to absorb --- clustering on
the customer identifier would produce clusters of size one and return the iid answer. There is no
time index, so there is nothing for a HAC correction to correct. What remains is the
cross-sectional worry, \emph{heteroskedasticity}: residual variance that differs across
customers, say because spend is more volatile among the loyal. The Huber--White sandwich with
\citeauthor{mackinnon1985}'s small-sample correction --- \code{HC1} --- is robust to exactly that,
and is the honest cross-sectional default \citep{white1980, mackinnon1985}.

Whether it \emph{matters} on this data is an empirical question, not a rhetorical one, so it is
tested. The Breusch--Pagan statistic \citep{breusch1979}, whose null is homoskedasticity, gives
$p = \nbFiveBpP$: no evidence against constant variance. And the sandwich changes the standard
error by a factor of $\nbFiveSeRatio$ --- a tenth of one per cent. On this data-generating
process the noise really is constant-variance Gaussian, and the robust covariance bought
\emph{nothing}.

It is kept anyway, and the reasoning is worth stating because it is the general case. One cannot
know in advance that the errors are homoskedastic without checking; HC1 costs nothing when they
are; and it is the only thing standing between the analyst and a badly understated interval when
they are not. Insurance that turns out to have been unnecessary was still worth buying. What must
\emph{not} happen is the reverse --- reporting a robust standard error and declining to say
whether it moved anything.

\subsection*{Two fits}

\Cref{tab:ctrl:sets} reports both, along with three further adjustment sets that
\cref{sec:ctrl:wrong} will need. Read the first two rows.

\begin{itemize}[leftmargin=*, itemsep=.2em]
  \item $Z = \varnothing$, the \textbf{naive} read --- spend on email and nothing else, the
        number a dashboard hands you: \eur{\nbFiveNaiveEst}, 90\,\% interval
        $[\nbFiveNaiveLo,\ \nbFiveNaiveHi]$, standard error \eur{\nbFiveNaiveSe}.
  \item $Z = \{\text{loyalty}\}$, the \textbf{backdoor-adjusted} read --- the set
        \cref{prop:ctrl:backdoor} endorses: \eur{\nbFiveAdjEst}, 90\,\% interval
        $[\nbFiveAdjLo,\ \nbFiveAdjHi]$, standard error \eur{\nbFiveAdjSe}.
\end{itemize}

Both are OLS. Both are competently executed. Both carry an honest, heteroskedasticity-robust
interval. One lands on the planted \eur{\nbFiveTrueAte} and covers it; the other misses it by
$\nbFiveNaiveErr$ and \emph{excludes} it. And the estimator has no idea which is which.

Look at \emph{how} the naive fit fails, because the manner of the failure is the chapter's
thesis. It is not vague. It is not hedged. Its interval is \eur{\nbFiveNaiveWidth} wide against
the adjusted fit's \eur{\nbFiveAdjWidth} --- the same order --- and it sits three and a half
euros from the truth, which it rules out. That is the definition of \textbf{confidently wrong}.
The bias from omitting \code{loyalty} is a fixed offset, and a confidence interval quantifies
sampling noise around whatever the estimator converges to --- never the distance between that
limit and the truth. \textbf{Precision is not validity}, and no standard error, however robust,
will tell you which of these two numbers to believe.

\begin{proposition}[Omitted-variable bias, binary treatment]\label{prop:ctrl:ovb}
Under \cref{eq:ctrl:outcome}, with $G$ independent of $T$, the naive coefficient converges to
%
\begin{equation}
  \hat\tau_{\text{naive}} \;\xrightarrow{\ p\ }\;
  \tau \;+\; \beta_L\,\Big(\E[L \mid T{=}1] \;-\; \E[L \mid T{=}0]\Big),
  \qquad \beta_L = \nbFiveOvbBetaL ,
  \label{eq:ctrl:ovb}
\end{equation}
%
that is: \textbf{bias = (effect of the omitted variable on the outcome) $\times$ (its imbalance
across arms)}.
\end{proposition}

Both factors are positive here. Loyalty raises spend ($\beta_L = \nbFiveOvbBetaL > 0$ in
\cref{eq:ctrl:outcome}), and loyal customers are emailed more
($\sigma(0.8L)$ is increasing in $L$), so the graph \emph{plus the signs of its edges} predicts
the direction of the naive bias before any data exists: \emph{biased up}. The
data-generating coefficients turn that direction into a magnitude, and \cref{sec:ctrl:valid}
checks the arithmetic against the fit.

Finally, the boundary that recurs in every chapter of this book: a 90\,\% confidence interval is
a property of the \emph{procedure}, not a probability about $\tau$. The decision rule this
campaign runs on --- $\Prob(\tau > \text{cost}) \ge 0.9$ --- is a statement about the effect
itself, and it has no frequentist translation. It needs a posterior. That, and not a better point
estimate, is what the next section is for.

% =====================================================================================
\section{The Bayesian model}
\label{sec:ctrl:bayes}

The Bayesian arm targets the \emph{same} estimand, \cref{eq:ctrl:estimand}, under the \emph{same}
adjustment set, on the \emph{same} \nbFiveN{} rows. It is \cref{eq:ctrl:ols} with a likelihood
and priors attached:
%
\begin{align}
  \text{spend}_i \mid \beta, \sigma
    &\;\sim\; \mathcal N\big(\beta_0 + \tau\,\text{email}_i + \gamma\,\text{loyalty}_i,\ \sigma^2\big),
    \label{eq:ctrl:lik}\\
  \beta_0 &\;\sim\; \mathcal N(20,\ 20^2), \qquad
  (\tau, \gamma) \;\sim\; \mathcal N(0,\ 10^2), \label{eq:ctrl:prior_b}\\
  \sigma &\;\sim\; \text{HalfNormal}(10). \label{eq:ctrl:prior_s}
\end{align}
%
The priors are weak and tied to the units of the data: spend has a standard deviation of
\eur{\nbFiveSpendSd}, so a $\mathcal N(0, 10^2)$ prior on a per-email effect spans every
commercially conceivable value and several inconceivable ones, and a HalfNormal$(10)$ on the
residual scale is wide enough for any plausible noise level. Nothing here is doing any work, and
that is deliberate. \textbf{The prior is not the interesting object in this chapter; the
adjustment set is.} Inference is by the No-U-Turn sampler \citep{hoffman2014} in PyMC
\citep{salvatier2016}: \nbFiveNDraws{} post-warmup draws, $\hat R = \nbFiveRhat$, minimum bulk
effective sample size \nbFiveEss{} \citep{vehtari2021}, \nbFiveDivergences{} divergent
transitions.

The posterior mean is \eur{\nbFiveBayesEst} with a 94\,\% highest-density interval of
$[\nbFiveBayesLo,\ \nbFiveBayesHi]$.

Then the same machinery is pointed at the \textbf{kitchen sink}: the identical likelihood, the
identical priors and the identical sampler, on spend regressed against the email, loyalty,
\code{responded} and \code{opened\_email} --- the ``control for everything'' instinct, implemented
faithfully. It converges just as well ($\hat R = \nbFiveKsRhat$, minimum ESS
\nbFiveKsEss, \nbFiveKsDivergences{} divergences). Its posterior mean is \eur{\nbFiveKsBar}. The
truth falls \emph{outside} its 94\,\% HDI. \Cref{sec:ctrl:compare} reads the consequences.

% =====================================================================================
\section{Diagnostics and validation}
\label{sec:ctrl:valid}

\subsection*{The bias formula, verified}

\Cref{prop:ctrl:ovb} predicted the naive bias from the graph and the data-generating
coefficients, without refitting anything. The loyalty imbalance across arms in this sample is
$\E[L \mid T{=}1] - \E[L \mid T{=}0] = \nbFiveOvbImbalance$, so \cref{eq:ctrl:ovb} predicts a bias
of $\nbFiveOvbBetaL \times \nbFiveOvbImbalance = \nbFiveOvbPredicted$. The fitted gap between the
naive coefficient and the truth is $\nbFiveOvbObserved$. They agree to \eur{\nbFiveOvbAgree}.

That is the standard to hold an observational read to, and it is worth stating as a rule:
\textbf{if you cannot pre-state the sign of the bias a proposed control would create or remove,
you are not ready to defend the estimate.} And when the pre-stated sign and the data disagree,
suspect the graph --- which is the next diagnostic.

\subsection*{Falsifying the graph against the data}

\Cref{as:ctrl:graph} is the load-bearing one, and this is where it is put at risk. The declared
DAG green-lit $\{\text{loyalty}, \text{opened\_email}\}$: \code{opened\_email} is not a collider
\emph{in the graph as drawn}, because the graph as drawn does not contain $G$. And yet
\cref{tab:ctrl:sets} shows that exact set is biased by $\nbFiveOpenBarErr$. The software was not
wrong. It reasoned faithfully about the graph it was given, and the graph was incomplete.

The escape is that a DAG is not merely an assumption --- it is a \emph{falsifiable} one. Every
graph implies a list of conditional independences that must hold in the data. The declared graph
implies, among others,
%
\begin{equation}
  \text{opened\_email} \;\perp\!\!\!\perp\; \text{spend} \;\;\Big|\;\;
  \{\text{email},\ \text{loyalty}\} ,
  \label{eq:ctrl:implication}
\end{equation}
%
because in the graph as drawn, every route onward from the opened-email node passes through the
email node. Test it. Of the \nbFiveFalsifNTests{} implications the declared graph makes,
\nbFiveFalsifNViolations{} is rejected at the 5\,\% level, and it is \cref{eq:ctrl:implication},
at $p = \nbFiveFalsifP$. The data says the graph is wrong, and it says so without anyone knowing
that $G$ exists.

This is the most valuable diagnostic in the chapter, so it is worth being explicit about what it
did and did not do. It did not tell anyone \emph{what} the missing structure is --- a latent
engagement trait feeding both \code{opened\_email} and \code{spend}. It told them the declared
graph cannot be the true one, which is enough to refuse the adjustment set the graph endorsed. A
causal DAG is an assumption; let the machine choose the adjustment set from the graph, and then
\emph{falsify the graph against the data} before trusting the verdict.

\subsection*{Bias or noise? Repeated sampling settles it}

One sample cannot distinguish a modest bias from bad luck. \Cref{fig:ctrl:bars} refits all five
adjustment sets on \nbFiveNSeeds{} fresh samples from \cref{eq:ctrl:outcome}. The endorsed set is
centred on the truth --- its mean error is $\nbFiveSeedBiasDag$, which is sampling noise. The
naive set sits $\nbFiveSeedBiasNaive$ away, sample after sample, and the kitchen sink
$\nbFiveSeedBiasKs$. Those offsets do not shrink and do not wander: they are properties of the
identification, not of the draw.

\input{build/floats/nb05_bars}

\subsection*{``But we have millions of rows''}

The claim that a wrong adjustment set is not repaired by data is the one every vendor deck
implicitly denies, so it is made formal. For a control set $W$, the OLS estimator obeys
%
\begin{equation}
  \hat\tau_W \;\xrightarrow{\ p\ }\; \tau \;+\; b_W ,
  \label{eq:ctrl:asymptotic}
\end{equation}
%
where $b_W$ is a \emph{fixed asymptotic bias} determined entirely by identification --- which
backdoor paths $W$ leaves open, and which collider paths it opens. As $n \to \infty$ the sampling
noise vanishes and $b_W$ does not move an inch.

\Cref{fig:ctrl:nsweep} sweeps the sample size to verify it. At $n = \nbFiveNSmall$ the naive
error is $\nbFiveSweepNaiveSmall \pm \nbFiveSweepNaiveSmallSd$; at $n = \nbFiveNLarge$ it is
$\nbFiveSweepNaiveLarge \pm \nbFiveSweepNaiveLargeSd$. The kitchen sink goes from
$\nbFiveSweepKsSmall$ to $\nbFiveSweepKsLarge$. The endorsed set goes from
$\nbFiveSweepDagSmall$ to $\nbFiveSweepDagLarge$. Every error bar narrows --- but only the
endorsed set narrows \emph{around zero}. The other two become more \emph{confidently} wrong.

\input{build/floats/nb05_nsweep}

\textbf{More data buys precision, never validity.} The first question to ask of any vendor's
attribution deck is therefore not \emph{how many rows} but \emph{what did you control for, and
what did that choice condition on?}

% =====================================================================================
\section{Point estimate versus posterior}
\label{sec:ctrl:compare}

\input{build/tables/nb05_arms}

\subsection*{On the endorsed set, they agree --- and it is worth saying plainly}

\Cref{tab:ctrl:arms} puts the classical fit and the posterior on the same estimand, the same
adjustment set and the same rows, with the classical interval taken at the posterior's 94\,\%
level so that the widths are comparable. The two point estimates differ by \eur{\nbFiveArmGap}
--- $\nbFiveArmGapSes$ classical standard errors --- and the interval widths differ by
\eur{\nbFiveArmWidthGap}. The truth is inside both.

That is not a disappointment; it is the point. With \nbFiveN{} rows and priors this weak, the
posterior \emph{is} the likelihood and the likelihood \emph{is} the least-squares fit --- the
Bernstein--von Mises phenomenon \citep{gelman2013}, showing up here in its most ordinary form.
\textbf{The Bayesian layer did not move the number, because the number was never where the
difficulty lay.} Anyone who arrived hoping that ``going Bayesian'' would improve the estimate has
just been shown, on their own data, that it does not.

\subsection*{What the posterior does buy}

Precisely one thing, and it is the thing the business asked for: $\Prob(\tau > \text{cost})$. It
is \nbFivePCost{} --- a probability \emph{about the effect}, natively a posterior quantity.
The classical arm cannot produce it. Not with a tighter interval, not with a bootstrap (which
approximates the sampling distribution of the \emph{estimator}, not a distribution of belief over
$\tau$), not with a $p$-value. The cell is not small; it is \emph{undefined}, and the decision
rule of \cref{sec:ctrl:euro} consumes nothing else.

Be equally precise about what does \emph{not} belong in that column. The sensitivity contour and
the E-value of \cref{sec:ctrl:euro} are arithmetic on the observed point estimate and would run
identically off the classical fit. They are sensitivity analysis, not Bayesian inference.
Crediting them to the posterior would be padding the bill.

\subsection*{What Bayes did not buy --- the chapter's thesis, demonstrated}

The third row of \cref{tab:ctrl:arms} is the same machinery pointed at the kitchen-sink
adjustment set. Read it slowly, because every column is a refutation of something practitioners
believe.

\begin{enumerate}[leftmargin=*, itemsep=.3em]
  \item \textbf{It converged immaculately.} $\hat R = \nbFiveKsRhat$, minimum ESS
        \nbFiveKsEss{}, \nbFiveKsDivergences{} divergences. There is nothing wrong with the
        \emph{inference}. The inference is faithfully reporting the posterior of a quantity that
        is not the causal effect.
  \item \textbf{It is centred on the wrong number.} Posterior mean \eur{\nbFiveKsBayesEst},
        an error of $\nbFiveKsBayesErr$, and the planted \eur{\nbFiveTrueAte} lies
        \emph{outside} its 94\,\% HDI $[\nbFiveKsBayesLo,\ \nbFiveKsBayesHi]$.
  \item \textbf{It is not punished with a wider interval.} This is the sentence to take away. Its
        HDI is \eur{\nbFiveKsBayesWidth} wide against the correct model's
        \eur{\nbFiveBayesWidth}. The extra six cents is not a warning; it is nothing. A model
        that is confidently wrong is \emph{not} flagged by its own uncertainty, because
        uncertainty prices sampling noise and this is not sampling noise
        (\cref{eq:ctrl:asymptotic}).
  \item \textbf{It emits a perfectly well-formed decision probability --- pointing the other
        way.} $\Prob(\tau > \text{cost}) = \nbFivePCostKs$ against the correct model's
        \nbFivePCost{}: \nbFiveDecisionKs{} against \nbFiveDecision. Same machinery, same data,
        same rule, opposite call.
\end{enumerate}

\input{build/floats/nb05_posterior}

\Cref{fig:ctrl:posterior} draws the two posteriors on one axis, and it is the picture the chapter
exists to produce: two tight, well-behaved, indistinguishable-looking distributions, one of which
is the causal effect and one of which is not.

A collider is a collider under any inference paradigm. Priors do not close backdoor paths;
posteriors do not detect open ones; MCMC converges just as happily on a biased estimand as on an
unbiased one, and reports its uncertainty just as confidently. \textbf{Bayes buys a decision
quantity; it does not buy identification.} Identification came from the graph in
\cref{sec:ctrl:ident}, and when the graph is wrong --- as \cref{sec:ctrl:valid} showed it can be
--- the posterior inherits the error and dresses it in a credible interval.

This is the same lesson the synthetic-control chapter teaches from the other side, and the pair is
worth holding together. There, the Bayesian arm's \emph{interval} failed because the likelihood was
misspecified, while its centre was fine. Here, the interval is fine and the \emph{centre} is
wrong, because the estimand was misspecified. Neither is a defect of the paradigm. Both are
defects of what was asserted before the sampler started.

% =====================================================================================
\section{The decision, in euros}
\label{sec:ctrl:euro}

\subsection*{The phantom revenue}

The campaign goes to \nbFiveCampaignN{} customers. The naive read says each email is worth
\eur{\nbFiveNaiveEst}; the backdoor-adjusted read says \eur{\nbFiveAdjEst}. The difference,
\eur{\nbFivePhantomPerCustomer} per customer, is not a rounding error at scale:
%
\begin{equation}
  \big(\hat\tau_{\text{naive}} - \hat\tau_{\{\text{loyalty}\}}\big) \times \nbFiveCampaignN
  \;=\; \eur{\nbFivePhantom} .
  \label{eq:ctrl:phantom}
\end{equation}
%
That is \eur{\nbFivePhantom} of incremental revenue booked into a plan, defended in a review, and
capable of justifying a budget --- and it does not exist. It is the confounding, sold as lift.
The cost of skipping the graph is not epistemic. It is that number.

\subsection*{The break-even, and the sets that flip it}

At a fully loaded \eur{\nbFiveCost} per contact, the decision rule is the one this book uses
throughout,
%
\begin{equation}
  \Prob(\tau > c) \;\ge\; 0.9 \;\;\Rightarrow\;\; \textsc{go},
  \qquad c = \eur{\nbFiveCost},
  \label{eq:ctrl:rule}
\end{equation}
%
and under the endorsed adjustment set the posterior gives
$\Prob(\tau > 0) = \nbFivePPos$ and $\Prob(\tau > c) = \nbFivePCost$: \textbf{\nbFiveDecision}.
The entire 94\,\% HDI sits above the cost line, so this is not a knife-edge call --- had the HDI
straddled the line, the honest answer would have been \emph{test before you spend}.

Now the uncomfortable arithmetic. Of the \nbFiveNSets{} adjustment sets in
\cref{tab:ctrl:sets}, \nbFiveNBelowCost{} falls below the \eur{\nbFiveCost} break-even --- the
kitchen sink, at \eur{\nbFiveKsBar}. \textbf{Over-controlling drags a genuinely profitable
campaign below its cost line and kills it.} The control choice is not a technicality that
statisticians care about; it is the difference between running the campaign and not running it.

\subsection*{How wrong could an unobserved confounder make this?}

\Cref{as:ctrl:backdoor} is untestable, so the residual risk is \emph{bounded} rather than
dismissed. Two parameterizations of the same worry are reported, and they are genuinely two ---
one does not summarise the other, because they live on different scales.

The first is additive and in euros. A hidden confounder $U$ with treatment-association $\gamma$
and outcome-association $\delta$ shifts the estimate to
$\text{observed} - \gamma\delta$; sweeping the pair over $(0, \nbFiveGammaMax)^2$ maps the region
in which the adjusted effect crosses zero (\cref{fig:ctrl:sensitivity}, left). Since the observed
effect is \eur{\nbFiveObservedAte}, that requires $\gamma\delta \approx \nbFiveObservedAte$ ---
a hidden variable more strongly associated with both email and spend than loyalty itself is.

The second is \citeauthor{vanderweele2017}'s \textbf{E-value} \citep{vanderweele2017}, which
lives on the risk-ratio scale and therefore requires the euro effect to be standardized first:
$d = |\hat\tau - c| / \sd(\text{spend}) = \nbFiveDStd$ outcome standard deviations, mapped to a
risk ratio by $\text{RR} \approx \exp(0.91\,d)$. The E-value is the minimum association --- with
\emph{both} treatment and outcome --- that an unmeasured confounder would need in order to move
the estimate to the threshold. Against the null it is \nbFiveEValue. Against the
\eur{\nbFiveCost} break-even it is \nbFiveEValueBe.

\input{build/floats/nb05_sensitivity}

The gap between those two numbers is the finding. It takes a very strong hidden confounder to
make this effect disappear; it takes a fairly ordinary one to make it \emph{unprofitable}.
``Robust to confounding'' is therefore not a property of an estimate. It is a property of an
estimate \emph{relative to a threshold}, and the threshold that matters is the price, not zero. Related sensitivity machinery on the regression scale ---
partial $R^2$ bounds benchmarked against an observed covariate --- is developed by
\citet{cinelli2020} and is the natural next instrument when the covariates are richer than they
are here.

A high E-value means \emph{robust to plausible hidden confounding}. It never means
\emph{unconfounded}. It is paired with the domain expert's judgement about what could be
missing, and with nothing else.

% =====================================================================================
\section{What can go wrong}
\label{sec:ctrl:wrong}

\input{build/tables/nb05_sets}

\Cref{tab:ctrl:sets} is the chapter in one table: the same OLS, the same HC1 covariance, the same
\nbFiveN{} rows, five adjustment sets. \nbFiveNCover{} of the \nbFiveNSets{} intervals cover the
truth. The widths span only \eur{\nbFiveWidthMin} to \eur{\nbFiveWidthMax} --- so the four wrong
sets are not vaguer than the right one, not noisier, not flagged. They are precise, and they are
precise about a different number.

\subsection*{The mediator: adjusting away the effect you are measuring}

In the mediator world the email is randomized, so there is no confounding whatsoever and the
no-controls regression is unbiased for the total effect --- and lands on it,
\eur{\nbFiveMedNone} against a truth of \eur{\nbFiveMedTotalTrue} (\cref{tab:ctrl:minis}).
Helpfully adding \code{visited} to the model does not improve precision. It \textbf{changes the
estimand}: the coefficient now answers \emph{what would the email do if site visits could be held
frozen?} --- the direct effect, \eur{\nbFiveMedCtrl} --- and silently discards the
\eur{\nbFiveMedIndirectTrue} of lift that flows through the visit.

The sin is not computing a direct effect. The funnel-mediation chapter builds the
total/direct/indirect decomposition on purpose, and it is often exactly what a marketer wants
\citep{imai2010}. The sin is computing one
\emph{by accident} and calling it the campaign's lift. An analyst who controls for post-send
engagement books only the fraction of the lift that bypasses the site.

\input{build/tables/nb05_minis}

\subsection*{M-bias: the pre-treatment variable that is still a trap}

Now the case that kills the most common heuristic in applied work. Last year's coupon-redemption
count is measured \emph{before} the send. It correlates with treatment ($\nbFiveMbCorrT$) and
with the outcome ($\nbFiveMbCorrY$). By every folk test it is a confounder, and by the folk rule
--- \emph{control for everything measured before treatment} --- it goes into the model.

It is in fact a collider between two latent traits: deal-proneness $U_1$, which drives both
coupon history and who the targeting model emails, and brand affinity $U_2$, which drives both
coupon history and spend. The only $T$--$Y$ path,
$T \leftarrow U_1 \to W \leftarrow U_2 \to Y$, contains a collider at $W$ and is therefore
\emph{already blocked}. The naive estimate is unbiased before anyone touches anything:
\eur{\nbFiveMbNone} against a truth of \eur{6}. Conditioning on $W$ opens the path --- once $W$
is known, a high $U_1$ implies a low $U_2$, so the two latent causes become correlated, splicing
$T \leftarrow U_1$ onto $U_2 \to Y$ --- and the estimate moves to \eur{\nbFiveMbCtrl}, a bias of
$\nbFiveMbBias$ \citep{ding2015}.

Two honest caveats, before anyone starts deleting covariates. First, \textbf{magnitude}: pure
M-bias is typically modest next to confounder bias --- compare the sub-euro shift here with the
$\nbFiveNaiveErr$ of confounding bias in \cref{tab:ctrl:sets}, and see \citet{ding2015} for the
general result. Second, \textbf{the common tangle}: when a variable is plausibly \emph{both} a
genuine confounder and an M-collider, adjusting is usually the lesser evil. The point is not
``never control for pre-treatment variables''. It is that \textbf{the decision comes from the
graph, not from the timestamp}. A timestamp is a heuristic; the graph is the criterion.

\subsection*{The collider nobody types into a regression: the filter}

Nobody on a marketing team puts \code{responded} into a model and feels clever about it. What
teams do instead, constantly, is \emph{filter}: \emph{let us measure the email's effect among the
customers who engaged.} The two acts are identical. \textbf{Subsetting a dataset is conditioning
on the subsetting variable} --- and it is worse than the regression version, because no model
summary will ever list \code{responded} as a covariate. The bias is invisible in the output.

\Cref{tab:ctrl:filter} runs it. Keeping only the \pc{\nbFiveSubPct} of customers who responded ---
\nbFiveSubN{} of \nbFiveN{}, hardly a niche slice --- conditions every downstream estimate on the
\code{responded} collider, and it does so silently. The naive read inside the subset is
\eur{\nbFiveSubNaiveEst}. And now the line that matters: handing the filtered analysis the
\emph{exactly correct} confounder still leaves a bias. It returns \eur{\nbFiveSubAdjEst},
$\nbFiveSubAdjErr$ off, with an honest HC1 interval that \emph{still covers the truth}. In a
single sample, nothing looks wrong.

\input{build/tables/nb05_filter}

It is the fresh-sample sweep that convicts it. Across \nbFiveSubNs{} new samples the shortfall
reappears every time: mean bias $\nbFiveSubMeanBias$ (sd \eur{\nbFiveSubSdBias}), which is
\nbFiveSubT{} standard errors from zero. That is the signature that separates a bias from noise,
and it is the same signature \cref{fig:ctrl:bars} draws for the regression version.

The pattern is worth recognising in the wild, because it rarely announces itself:

\begin{itemize}[leftmargin=*, itemsep=.15em]
  \item \textbf{converter-only attribution} --- ``of the customers who purchased, which
        touchpoints did they see?'';
  \item \textbf{opened-email lift reads} --- ``the subject line's effect among opens'';
  \item \textbf{survey-based incrementality} --- only responders answer the survey;
  \item \textbf{``active users'' dashboards} --- activity is caused by both the campaign and the
        outcome it proxies.
\end{itemize}

The fix is never a better covariate \emph{inside} the subset. It is refusing the filter ---
estimating on everyone the campaign could have touched --- or modelling the selection explicitly.

\subsection*{The kitchen sink}

Which brings the chapter back to the instinct it opened with. ``Everything'' --- loyalty,
\code{responded} and \code{opened\_email} together --- is the worst of the five sets:
\eur{\nbFiveKsBar}, an error of $\nbFiveKsBarErr$, and a 90\,\% interval of
$[\nbFiveKsBarLo,\ \nbFiveKsBarHi]$ that excludes the truth. It stacks the collider bias from
\code{responded} on top of the latent-collider bias from \code{opened\_email}. It is the only set
that flips the campaign decision. And it is what a well-meaning analyst produces when told to be
thorough.

\subsection*{A practitioner's covariate checklist}

The toy variables map one-to-one onto the columns of a real campaign table.

\begin{center}\small
\begin{tabular}{@{}p{0.42\linewidth} p{0.30\linewidth} p{0.20\linewidth}@{}}
\toprule
\textbf{Variable, as measured} & \textbf{Likely causal role} & \textbf{Action} \\
\midrule
Loyalty tier, RFM score computed \emph{before} the send, region, seasonality, prior-quarter spend
  & Confounder: drives both targeting and spend & \textbf{control} \\[2pt]
Opens, clicks, site visits \emph{after} the send
  & Mediator --- or a hidden collider via latent engagement & \textbf{exclude} when the question
    is the total effect \\[2pt]
Responded, converted, unsubscribed
  & Collider (post-outcome) & \textbf{exclude}, and never \emph{filter} on it \\[2pt]
Survey participation, post-campaign app install
  & Selection collider & \textbf{exclude}; distrust any panel filtered on them \\[2pt]
Pre-treatment \emph{proxy} of an unmeasured trait (historical open rate, coupon history)
  & It depends --- draw the graph & M-bias risk if it is a common effect of latents \\
\bottomrule
\end{tabular}
\end{center}

\medskip\noindent
Pocket rule: \textbf{pre-treatment \emph{and} a plausible common cause of $T$ and $Y$ $\Rightarrow$
in; caused by $T$ $\Rightarrow$ out; anything else $\Rightarrow$ draw the DAG.}

% =====================================================================================
\section{On real data}
\label{sec:ctrl:real}

Everything above was graded against a truth that was planted. The objection writes itself: a
method that recovers an effect somebody put there is not thereby proved to work. There is one
dataset that answers it, and it has been humbling the profession since 1986.

\citet{lalonde1986} took a \emph{randomized} evaluation of a job-training programme --- whose
experimental effect on 1978 earnings is roughly \usd{\nbFiveLlBenchmark} --- threw away the
experimental control group, and replaced it with a comparison group drawn from population
surveys. The observational data that results is a trap with a known answer: any adjustment
strategy can be \emph{scored} against the experiment. \citet{dehejia1999} made the composite the
standard benchmark it remains. The version used here has \nbFiveLlN{} observations:
\nbFiveLlNTreated{} randomized trainees against \nbFiveLlNControl{} survey comparisons.

\Cref{tab:ctrl:lalonde} reads the same observational data three ways.

\input{build/tables/nb05_lalonde}

\begin{description}[leftmargin=1.1em, style=nextline, font=\normalfont\itshape]
  \item[The naive comparison is negative.] \usd{\nbFiveLlNaive}. Not merely wrong in magnitude ---
  wrong in \emph{sign}. Read literally, job training \emph{cut} earnings. The reason is printed in
  the same table's caption: in 1974, years before the programme existed, the future trainees
  earned \usd{\nbFiveLlPreTreated} and the comparison group \usd{\nbFiveLlPreControl}. The groups
  were never comparable, and the raw gap is measuring who was poor, not what training did.
  \item[The responsible-looking adjustment lands short.] Six pre-treatment demographic
  covariates --- age, education, race, marital status, degree --- repair the sign and give
  \usd{\nbFiveLlBad}, still \usd{\nbFiveLlGapBad} below the experimental benchmark. It looks
  diligent. It is missing the load-bearing variables.
  \item[The backdoor-adjusted read lands.] Restoring the two pre-programme earnings years
  (\code{re74}, \code{re75}) --- \nbFiveLlNGood{} covariates in total --- gives \usd{\nbFiveLlGood},
  within \usd{\nbFiveLlGapGood} of the randomized answer.
\end{description}

Same data, three defensible-\emph{sounding} specifications, three different stories. And choosing
among them is exactly the skill of \cref{sec:ctrl:ident}, exercised without a simulator: prior
earnings are the strongest plausible common cause of \emph{enrolling in a training programme} and
of \emph{earning later}, so \cref{prop:ctrl:backdoor} demands them in $Z$. The demographic-only
specification leaves that backdoor wide open.

\begin{remark}[The mirror image]
Dropping a load-bearing confounder and adding a collider are the same kind of error --- an
identification failure --- and neither is repaired by a better estimator or a larger sample
(\cref{eq:ctrl:asymptotic}). The simulated chapter demonstrates the second; LaLonde demonstrates
the first, on data that has resisted forty years of increasingly sophisticated estimators. The
literature's long detour through matching, weighting and machine-learning adjustments on this
dataset is instructive precisely because the estimators changed and the \emph{answer} tracked the
\emph{covariate set}.
\end{remark}

% =====================================================================================
\section{Takeaways}
\label{sec:ctrl:take}

\begin{enumerate}[leftmargin=*, itemsep=.45em]
  \item \textbf{``Control for everything you have'' is wrong, and it is expensive.} The kitchen
        sink returns \eur{\nbFiveKsBar} against a true \eur{\nbFiveTrueAte}, excludes the truth
        from its interval, and \emph{reverses} the campaign decision
        (\nbFiveDecisionKs{} against \nbFiveDecision). The instinct that feels like caution is
        the one that kills a profitable campaign.

  \item \textbf{The role decides, not the timestamp.} Confounders must be adjusted for; mediators,
        colliders and descendants of treatment must not. A collider path is \emph{closed until you
        condition on it}. And a variable measured \emph{before} treatment can still be a collider:
        the M-bias example moves a correct estimate by $\nbFiveMbBias$ using a covariate that
        correlates $\nbFiveMbCorrT$ with treatment and $\nbFiveMbCorrY$ with the outcome and
        passes every folk test for a confounder.

  \item \textbf{Filtering is conditioning.} Restricting to the \pc{\nbFiveSubPct} of customers who
        responded opens a collider path that no covariate inside the subset can close. Even given
        the correct confounder, the filtered fit is biased by $\nbFiveSubMeanBias$ across
        \nbFiveSubNs{} samples --- \nbFiveSubT{} standard errors from zero. Converter-only
        attribution, opened-email lift reads and ``active user'' dashboards are all the same
        error in different clothes.

  \item \textbf{No diagnostic in the output will tell you the adjustment set is wrong.} The wrong
        models here converge perfectly ($\hat R = \nbFiveKsRhat$, \nbFiveKsDivergences{}
        divergences), fit well, and produce intervals of the same width as the right one --- the
        five 90\,\% intervals span \eur{\nbFiveWidthMin} to \eur{\nbFiveWidthMax}, and only
        \nbFiveNCover{} of \nbFiveNSets{} covers the truth. \textbf{A confidently wrong model is
        not punished with a wide interval}, because an interval prices sampling noise and a
        wrong adjustment set is a fixed offset.

  \item \textbf{More data buys precision, never validity.} $\hat\tau_W \to \tau + b_W$, and $b_W$
        is set by identification alone. At $n = \nbFiveNLarge$ the naive read is still
        $\nbFiveSweepNaiveLarge$ out and the kitchen-sink read $\nbFiveSweepKsLarge$ out --- more
        confidently wrong than at $n = \nbFiveNSmall$. The first question for a vendor deck is not
        how many rows.

  \item \textbf{Bayes buys a decision quantity, not identification.} On the endorsed set the two
        arms agree to \eur{\nbFiveArmGap} and their intervals differ by \eur{\nbFiveArmWidthGap}:
        the sampler adds no causal content, and the chapter says so without flinching. What it
        adds is $\Prob(\tau > \text{cost}) = \nbFivePCost$, which the classical arm cannot produce
        at any level of effort --- and which the \emph{same} sampler reports just as confidently,
        as \nbFivePCostKs, for an estimand that is not the causal effect.

  \item \textbf{The DAG is an assumption --- so falsify it.} The graph declared here was wrong: it
        could not see the latent trait behind \code{opened\_email}, and it green-lit an adjustment
        set that is biased by $\nbFiveOpenBarErr$. The data caught it anyway.
        \nbFiveFalsifNViolations{} of the \nbFiveFalsifNTests{} conditional independences the
        graph implies is rejected, at $p = \nbFiveFalsifP$. Let the machine choose the adjustment
        set from the graph; then test the graph against the data before believing either.

  \item \textbf{Bound what you cannot test.} \Cref{as:ctrl:backdoor} is untestable, so it is
        bounded: an unmeasured confounder would need an association of \nbFiveEValue$\times$ with
        both email and spend to null the effect --- but only \nbFiveEValueBe$\times$ to drag it
        below the \eur{\nbFiveCost} break-even. Robustness is relative to a threshold, and the
        threshold that matters is the price.

  \item \textbf{The euro anchor.} Budgeting on the naive read books \eur{\nbFivePhantom} of
        incremental revenue that does not exist across a \nbFiveCampaignN-customer campaign. That
        is what an unclosed backdoor costs, and it was produced by a decision made before any model
        ran.
\end{enumerate}

\notebooknote{%
  \Crefrange{sec:ctrl:dgm}{sec:ctrl:wrong} and \cref{sec:ctrl:real} are
  \code{notebooks/05\_what\_to\_control\_for.ipynb}, which runs in a fast teaching mode
  (\code{CMP\_FAST=1}) and a full mode (\code{CMP\_FAST=0}); every number in this chapter comes
  from the full run, emitted by \code{cmp.report} and injected here as a macro. \Cref{sec:ctrl:real}
  fetches the \citeauthor{dehejia1999} composite of \citeauthor{lalonde1986}'s data over the
  network and is gated on \code{CMP\_REAL=1}. The simulator is \code{cmp.dgp.dag\_control\_demo};
  the graph is declared to \code{pathmc}, which enumerates the admissible adjustment sets, flags
  colliders and tests the graph's implied conditional independences; the seed is fixed in the
  notebook.}
